\section{When counting fails}
\label{sec:counting}

Everything so far took the count for granted. A family is a set of members
with a numbering, and the count at a point $x$ is the number $\nu$ such that
$x$ lies on the interpolated member $\nu$: the solution of an equation
$F(x,\nu)=0$, ``$x$ is on member $\nu$''. Whether that equation has a
solution, and one solution, is not automatic, and the ways it can fail are
the whole taxonomy of families. There are exactly three.

\begin{lookup}{the implicit function theorem, and winding}
What to look up: \emph{the implicit function theorem}. The one consequence
used: wherever $\partial F/\partial\nu\ne0$, the equation $F(x,\nu)=0$
defines $\nu$ as a smooth function of $x$ near that point; where
$\partial F/\partial\nu=0$ it may not. And \emph{winding number}: walk once
round a point and the angle you make with it gains a full turn, $2\pi$,
which no continuous choice of angle can undo.
\end{lookup}

\begin{theorem}[Trichotomy]
\label{thm:trichotomy}
The count $\nu(x)$ of a family exists in exactly one of three ways.
\emph{As a formula:} a single smooth function on the whole plane (the exact
rung: lines, rings, rays, spirals, waves). \emph{Up to a whole number:} a
function on the plane with a point removed, which gains a whole number each
time a loop goes round the point (the winding rung: the polar angle is the
prototype, and a family built on it has a \emph{defect} at the point, with
a quantized charge). \emph{Only as a search:} where two members pass through
the same point at the same interpolated index, $\partial F/\partial\nu=0$,
the count branches, no relabelling repairs it, and the nearest member must
be found by looking (the fold rung: walking families, whose members are
each displaced and turned a little more than the last).
\end{theorem}

\begin{proof}
Where $\partial F/\partial\nu\ne0$ the lookup gives $\nu$ locally as a smooth
function of $x$. The local pieces glue into one function on the plane
exactly when carrying a value round any loop brings it back to itself; if
it comes back changed, the change is a whole number (the index of the same
member), it depends only on which points the loop encloses, and it is the
winding of the family about those points. Where $\partial F/\partial\nu=0$
two members meet $x$ at the same index to first order; the local pieces join
there instead of overlapping, and no cover of the plane makes the count a
function again.
\end{proof}

\paragraph{Defects.}
Add to a family's count the polar angle about a point divided by a full
turn, times a whole number $A$. Nothing about the family's members changes
away from the point, but the count now gains $A$ around any loop that
encloses it, and so does the difference $D$ against an unmodified twin. The
bands of that pair, the level lines of $D$, must therefore \emph{end}: $A$
of them inside any loop round the point, and, since winding adds, the
signed number of endings inside any loop is $A$ times the number of enclosed
points. That is the fork grating of optics~\cite{NyeBerry1974Dislocations}
built from a count rather than a wavefront, and the count says one more
thing: near the point the added count changes at a rate $A/2\pi r$, so the
heterodyne ratio of the beat is $\het=As/2\pi r$ for a family of pitch $s$,
and the bands cannot exist inside the core $r^\star=2As/\pi$, where the beat
is no slower than the carrier. Measured on \defProbes\ probes with charges
up to \defChargeMax: the count of endings is exact, and the core radius
holds to within $\defCoreWorst$ (script \texttt{defects.mjs};
Figure~\ref{fig:defects}).

\begin{figure}[t]
  \centering
  \includegraphics[width=0.32\linewidth]{defect-fork.png}\hfill
  \includegraphics[width=0.32\linewidth]{defect-pair.png}\hfill
  \includegraphics[width=0.32\linewidth]{defect-envelope.png}
  \caption{\figlead{A count that winds.} Left: a family of lines whose count
  has gained five around the centre, the fork. Middle: the same family over
  an unmodified twin. Right: what pooling sees of the pair, with five bands
  ending at the defect. The number of endings inside any loop is the
  winding it encloses, exactly. Script \texttt{defectfigs.mjs}.}
  \label{fig:defects}
\end{figure}

\paragraph{Folds.}
A walking family is a nest of curves --- circles, squares, hexagons --- each
grown by a step $s$, displaced by a step $\delta$ and turned by a step
$\theta$ more than the last. Its members are the fronts of something that
spreads at speed $s$ from a source that moves at speed $|\delta|$ per
member, and the family folds exactly where the front's normal speed
vanishes. Two consequences have closed forms. A front that only moves folds
if and only if the source is \emph{supersonic}: it moves faster than it
spreads, $\rho(-\delta)>s$ in the gauge's own metric, and the fold is the
shock along the wake. A front that only turns folds every shape but the
circle, at the radius $s/(\theta\sin(\pi/k))$ for the regular $k$-gon, where
a vertex of one member first crosses the next. Onsets measured by brute
force on families that never saw the calculus are within
\foldWorstOnsetPct\% of these radii (script \texttt{foldlaw.mjs};
Figure~\ref{fig:foldlaw}). Beyond the fold the count is multivalued, the
number of branches through a point grows with the distance from the centre,
and no formula exists: a rational turn $2\pi a/b$ splits the family into $b$
nests that can each be solved, and any other turn is provably beyond
algebra. Searching for the nearest member is the price of a family whose
count has folded, and the instrument of \S\ref{sec:instrument} pays it with
a certificate rather than a guess.

\begin{figure}[t]
  \centering
  \includegraphics[width=\linewidth]{fold-law.png}
  \caption{\figlead{Where a count folds.} Predicted fold curves (crimson)
  drawn over brute-force renders of walking families, the onset radius
  dashed. Left, a supersonic displacement: the wake. Centre and right,
  turning hexagons and a turning ellipse: the fold where a vertex of one
  member crosses the next. Script \texttt{foldlaw-figs.mjs}.}
  \label{fig:foldlaw}
\end{figure}

One thing a field cannot do closes the taxonomy. A \emph{field} added to a
count --- any function of position, the polar angle included --- shifts
$\nu$ by a function of $x$ and leaves $\partial F/\partial\nu$ alone. So a
field can move a family within its rung or lift it from the exact rung to
the winding rung by a whole number, and it can never fold one. Folding is a
property of the members' geometry; winding is a property of the count; and
the third pattern, being a function of counts, inherits the winding and
ends where the count does.

\begin{tryit}{a count that winds, and one that folds}
In \Moire, take a single family of lines and type the polar angle divided
by a full turn into its field. One line now ends at the centre, and a
second family of the same pitch laid over it shows one band ending there:
charge one. Multiply the field by three and three bands end. Then take a
family of hexagonal rings and give each member a small turn relative to the
last: the nest is clean near the centre and folds where the theorem says,
at radius $s/(\theta\sin(\pi/6))$, and the \emph{ratio} view goes blank
inside the defect's core.
\end{tryit}
